\section{Interfaces and data contracts}\label{sec:5}

Three communication patterns are used, each for the job it fits. \textbf{MQTT pub/sub} carries
readings, commands, residuals, fault events and strategies, because one producer feeds several
consumers and a dead producer should go silent rather than block. \textbf{REST} carries every call
to BuildSim, because each is a one-to-one request whose result the caller must confirm.
\textbf{HTTP} serves the dashboard page and its JSON. This section is written to be precise enough
to re-implement a component from.

\subsection{MQTT topics and payloads}\label{sec:5.1}

\texttt{\{zone\}} is one of the six zone identifiers \texttt{A125A}, \texttt{A125},
\texttt{A116A}, \texttt{A115}, \texttt{A124}, \texttt{A145}; \texttt{\{ahu\}} is
\texttt{AHU-A120}. The zone identifier is the bare room name and not BuildSim's key
\texttt{level0/A115}, because the key contains a slash, which would split an MQTT topic; the
BuildSim adapter is the only place that knows the key.

\begin{longtable}[]{@{}
  >{\raggedright\arraybackslash}p{(\columnwidth - 8\tabcolsep) * \real{0.2261}}
  >{\raggedright\arraybackslash}p{(\columnwidth - 8\tabcolsep) * \real{0.1500}}
  >{\raggedright\arraybackslash}p{(\columnwidth - 8\tabcolsep) * \real{0.1600}}
  >{\raggedright\arraybackslash}p{(\columnwidth - 8\tabcolsep) * \real{0.2100}}
  >{\raggedright\arraybackslash}p{(\columnwidth - 8\tabcolsep) * \real{0.2539}}@{}}
\toprule
\textbf{Topic} & \textbf{Publisher} & \textbf{Subscribers} & \textbf{Payload fields} &
\textbf{Semantics} \\
\midrule
\endhead
\texttt{readings/\{zone\}/co2} & \texttt{co2-sensor} & controller, monitor, supervisor,
dashboard &
\texttt{zone}, \texttt{ppm}, \texttt{ts}, \texttt{quality} & One per room per second.
\texttt{quality} is the sensor process's own judgement of the reading. The supervisor is a
subscriber because it sizes a compensation plan from measured air \\
\texttt{readings/\{ahu\}/\allowbreak fan-power} & \texttt{co2-sensor} & monitor, dashboard & \texttt{ahu},
\texttt{watts}, \texttt{ts}, \texttt{quality} & One fan, so one topic. The monitor's only
cross-room signal \\
\texttt{commands/\{zone\}/damper} & \texttt{dcv-controller}; \texttt{supervisor} for the probe &
actuator, monitor, supervisor & \texttt{zone}, \texttt{position}, \texttt{ts}, \texttt{reason},
\texttt{commandId} & \texttt{commandId} is what the residual is keyed to; \texttt{reason} is
human-readable so every action is auditable. The supervisor subscribes to its own topic because
it has no other damper feedback \\
\texttt{commands/\{ahu\}/\allowbreak fan-speed} & \texttt{supervisor} only & actuator, monitor,
dashboard & \texttt{ahu},
\texttt{speed}, \texttt{ts}, \texttt{reason}, \texttt{commandId} & The compensation lever. The
inner loop never writes this topic; the monitor subscribes so that its believed flow split tracks
the speed it was not told about directly \\
\texttt{residuals/\{zone\}} & \texttt{fault-monitor} & dashboard & \texttt{zone},
\texttt{residual}, \texttt{cusum}, \texttt{ts}, \texttt{commandId} & Published every cycle,
including in \texttt{WARMUP}, where no diagnosis is issued \\
\texttt{faults/\{zone\}} & \texttt{fault-monitor} & supervisor, dashboard & \texttt{zone},
\texttt{state}, \texttt{classification}, \texttt{evidence}, \texttt{ts}, \texttt{caseId} &
\texttt{state} is the monitor state; \texttt{classification} is one of \texttt{STUCK\_DAMPER},
\texttt{DAMPER\_STUCK\_OPEN}, \texttt{BLOCKED\_FILTER}, \texttt{SENSOR\_DRIFT} or
\texttt{UNCLASSIFIED} \\
\texttt{strategy/\{zone\}} & \texttt{supervisor} & controller, dashboard & \texttt{zone},
\texttt{mode}, \texttt{ts}, \texttt{caseId} & \texttt{DEGRADED} tells the inner loop to stop
commanding that room; \texttt{COMPENSATING} carries the healthy rooms' trim \\
\texttt{service/\{zone\}} & \texttt{dashboard}, or any operator tool & \texttt{fault-monitor} &
\texttt{zone}, \texttt{caseId}, \texttt{reason}, \texttt{ts} & The equipment behind a latched case
was repaired, which is the \emph{only} thing that clears one. Refused unless that exact case is open
and latched. \textbf{Deliberately not under \texttt{commands/}}: it reports something about the
world rather than instructing an actuator, and putting it there would give a second publisher
authority over the topic space \texttt{damper-actuator} owns \\
\bottomrule
\end{longtable}

Two example payloads, with units in the field names or stated, so no consumer has to guess:

\begin{footnotesize}
\begin{verbatim}
readings/A115/co2        { "zone": "A115", "ppm": 1020.4,
                           "ts": "2026-10-05T09:14:00Z", "quality": "GOOD" }

faults/A115              { "zone": "A115", "state": "CONFIRMED",
                           "classification": "STUCK_DAMPER",
                           "evidence": ["cusum 175 ppm > h 150 ppm for 2 cycles",
                                        "fan power flat at 137 W",
                                        "5 of 6 residuals quiet",
                                        "overnight floor 418 ppm"],
                           "ts": "2026-10-05T09:21:00Z", "caseId": "case-3" }
\end{verbatim}
\end{footnotesize}

\subsection{REST, against BuildSim}\label{sec:5.2}

BuildSim's schema is not mine, so this is the one place a translation layer exists.

Six calls, all one-to-one writes or reads whose result the caller must confirm, which is why each
is REST and not a message: the simulator writes sensor values and reads back the actuator state
that actually applied, \texttt{co2-sensor} reads its own sensor, \texttt{damper-actuator} writes
actuator state as the single source of truth, \texttt{supervisor} writes the work-order alert, and
\texttt{dashboard} repaints the room layers every 5 s idempotently. \textbf{Appendix \ref{sec:C.3}} gives
caller, direction, purpose and the reason for each.

\subsection{Journal (TimescaleDB)}\label{sec:5.3}

Six typed hypertables: \texttt{readings}, \texttt{fan\_power}, \texttt{commands},
\texttt{residuals}, \texttt{decisions} and \texttt{strategies}. Every row of all six carries two
provenance columns (\texttt{source}, which process wrote it, and \texttt{received\_at}, distinct
from the event's own observation time so that an arrival delay is visible rather than invisible),
and \texttt{commands} and \texttt{residuals} additionally carry \texttt{command\_id}. That last
column is the one that makes R-5 mean anything: \textbf{a residual that cannot be tied to the
command it was judged against is not evidence}. Window queries over the residual series for one
\texttt{caseId} are hand-written native SQL rather than an ORM, which is the one place the plan
allows raw SQL and the one place it is clearly right: the query is a time-bounded ordered window,
not entity retrieval.

\textbf{There is no retention or compression policy}, and no chunk interval is specified, so
TimescaleDB's seven-day default applies and nothing is ever dropped. For runs of the length this
project uses that is harmless; it is stated here because describing a policy that does not exist
would be the easier and less honest option.
